% ============================================================
% conclusiones.tex — Conclusiones del documento
% SemArq1_TeamBales.tex, después del último capítulo (P6) y antes de
% \printbibliography.
% ============================================================

\chapter*{Conclusiones}
\addcontentsline{toc}{chapter}{Conclusiones}

El análisis de las preguntas lleva a una misma idea: lo que distingue a una buena arquitectura no es el número de capas ni la forma del dibujo, sino que la lógica de negocio no dependa de los detalles técnicos. Cuando se consigue, el sistema se puede probar y cambiar por partes sin que cada modificación se propague a todo el código. Esa independencia, sin embargo, tiene un coste de construcción y mantenimiento \fuente{martin2017}{218}, por lo que la arquitectura adecuada depende del tamaño del sistema y de cuánto se espera que cambie.
